% Source: https://artofproblemsolving.com/community/c3h2042685_ball_bouncing_problem
% Stable identity: aops-topic-2042685-q01
Put the box at \(0\le x\le50\), \(0\le y\le10\) with \(M=(25,0)\). By symmetry assume the ball moves to the right with slope \(k>0\). Instead of reflecting the ball, reflect the box: the path becomes a straight ray from \(M\) in the plane tiled by copies of the box, the corners become the lattice points \((50i,10j)\), and each bounce corresponds to crossing a grid line.

Suppose the ray first reaches a corner at \((50i,10j)\) with \(i,j\ge1\). On the way it crosses the vertical lines \(x=50,\dots,50(i-1)\) and the horizontal lines \(y=10,\dots,10(j-1)\), giving \((i-1)+(j-1)\) bounces. So \(i+j=5\), and
\[
k=\dfrac{10j}{50i-25}=\dfrac{2j}{5(2i-1)}\in\left\{\dfrac{8}{5},\ \dfrac{2}{5},\ \dfrac{4}{25},\ \dfrac{2}{35}\right\}
\]
for \((i,j)=(1,4),(2,3),(3,2),(4,1)\). (In each case no earlier lattice corner lies on the segment, so the corner is reached for the first time.)

Among the choices only \(1.6=\dfrac{8}{5}\) appears. For \(k=1.6\) the ball hits the top at \(x=31.25\), the bottom at \(x=37.5\), the top at \(x=43.75\), and then the corner \(B=(50,0)\):

\begin{asy}
usepackage("amsmath");
size(7cm);
defaultpen(fontsize(10pt));
real sc=0.12;
pair A0=(0,0), B0=(50,0), C0=(50,10), D0=(0,10);
draw(scale(sc)*(A0--B0--C0--D0--cycle),linewidth(0.9));
pair M=(25,0), P1=(31.25,10), P2=(37.5,0), P3=(43.75,10);
draw(scale(sc)*(M--P1--P2--P3--B0),red+linewidth(0.9));
dot(scale(sc)*M); label("$M$",scale(sc)*M,S); label("$A$",scale(sc)*A0,SW); label("$B$",scale(sc)*B0,SE);
label("$50$",scale(sc)*(12.5,0),S); label("$10$",scale(sc)*(0,5),W);
label("\mbox{}",(min(currentpicture,user=true)+max(currentpicture,user=true))/2,black); // keeps the text after the figure black when LaTeX typesets the labels inline
\end{asy}

The answer is \(\boxed{\textbf{(D) }1.6}\).
